% ==============================================================================
% T0 Theory / FFGFT: Document 270 (Chapter Content, DE)
% Titel: Projektionskette T^4 -> T^0 -- Zwei Typen von Dimensionsreduktion
% Autor: Johann Pascher
% Datum: Juni 2026
% Referenzen: Dok. 147, 159, 183, 185, 206, 254, 258, 259, 268
% ==============================================================================

\section*{Abstract}
Die Reduktion der fundamentalen 4-Torus-Struktur $T^4$ der FFGFT auf niedrigere
Dimensionen wird als formale Kette $T^4 \to T^3 \to T^2 \to T^1 \to T^0$
untersucht. Die zentrale Aussage ist, dass diese Kette \emph{nicht} aus
gleichartigen Schritten besteht. Sie zerf\"allt in zwei qualitativ verschiedene
Operationen. Der erste Schritt ($D4\!\to\!D3$) ist keine verlustbehaftete
Projektion, sondern eine \emph{Dualit\"ats-Dekompaktifizierung}: die kompakte
Massedimension $S^1_m$ wird \"uber die Fundamentalrelation $\tilde{T}\cdot m = 1$
in die nichtkompakte Zeitrichtung abgerollt. Diese Operation ist im Modeninhalt
informationserhaltend; einzig die direkte Sichtbarkeit der Kompaktheit
\"uberlebt als diskrete Spektralsignatur, kodiert durch die Rekursion $\xi^n$.
Die weiteren Schritte ($D3\!\to\!D2\!\to\!D1\!\to\!D0$) sind echte geometrische
Projektionen r\"aumlicher Richtungen: streng verlustbehaftet, nicht reversibel,
durch die Rekursion \emph{nicht} kompensiert. Der Schritt $D3\!\to\!D2$ wird als
holographische Projektion identifiziert. Die topologische Identifikation
$t\sim t+R_m$ wird als physikalische Observable nachgewiesen ($T_\text{rev}$,
CHSH-Abweichung $\sim\xi$). Der Abw\"artsreduktion wird die \emph{Aufw\"arts}-
Operation gegen\"ubergestellt: die bijektive repr\"asentationelle \"Ubersetzung
$T^4 \leftrightarrow$ Hilbertraum/Matrix (Typ~III, Dok.~230/206/231), die
\emph{verlustfrei} ist -- gerade weil sie keine geometrischen Extra-Dimensionen
hinzuf\"ugt, sondern denselben 4D-Inhalt in einer h\"oherdimensionalen Sprache
ausdr\"uckt. Die Arbeit l\"ost damit die scheinbare Spannung zwischen
``Projektion verliert Information'' und ``kompaktifiziert / dekompaktifiziert''
auf und ordnet die geometrische Abw\"artsreduktion gegen die repr\"asentationelle
Aufw\"arts\"ubersetzung ein.



% ============================================================================
\section{Motivation: die Spannung in der Fragestellung}
% ============================================================================

Die FFGFT beschreibt die fundamentale Physik auf einem kompakten 4-Torus $T^4$
mit dem einzigen dimensionslosen Parameter $\xi = 4/30000$ und der zentralen
Relation $\tilde{T}\cdot m = 1$. Die beobachtbare Welt ist hingegen ein
dreidimensionaler Raum mit Zeitentwicklung, $\mathbb{R}^3 \times \mathbb{R}_t$.

Drei Beobachtungen zu diesem \"Ubergang stehen scheinbar im Widerspruch:
\begin{itemize}
  \item[(A)] \textbf{Verlustbild.} Jede Projektion $D{\to}D{-}1$ ist nicht
  injektiv, verliert also Information. Die Kette $T^4 \to \dots \to T^0$ ist
  nicht reversibel.
  \item[(B)] \textbf{Kompaktheitsbild.} $T^4$ ist vollst\"andig kompaktifiziert;
  die beobachtbare Raumzeit $\mathbb{R}^3 \times \mathbb{R}_t$ ist
  dekompaktifiziert.
  \item[(C)] \textbf{Aufw\"artsbild.} Es existiert bereits eine \emph{bijektive}
  \"Ubersetzung von $T^4$ in einen h\"oherdimensionalen Raum -- den Hilbertraum
  (Dok.~230) -- die \emph{verlustfrei} ist.
\end{itemize}

Diese drei Bilder lassen sich nicht naiv \"ubereinanderlegen: W\"are der erste
Schritt reiner Verlust (A), k\"onnte daraus keine \emph{dekompaktifizierte}
Zeitrichtung (B) entstehen, denn Vergessen erzeugt keine neue Dimension; und es
bliebe r\"atselhaft, wie eine Aufw\"artsbewegung (C) zugleich verlustfrei sein
kann, w\"ahrend die Abw\"artsbewegung (A) Information vernichtet.

Die Aufl\"osung: Es gibt nicht eine, sondern \emph{drei} dimensions\"andernde
Operationen von \emph{zwei Arten}. \emph{Geometrische} Operationen (Typ~I:
Masse-Zeit-Dualit\"at; Typ~II: r\"aumliche Projektion) bauen die beobachtbare
Welt aus $T^4$ ab; die \emph{repr\"asentationelle} Operation (Typ~III)
\"ubersetzt $T^4$ verlustfrei nach oben. Die geometrisch-runter verlustbehaftete
gegen die repr\"asentationell-hoch bijektive Bewegung -- diese Asymmetrie ist das
zentrale Resultat der Arbeit.

% ============================================================================
\section{Struktur von $T^4$}
% ============================================================================

\subsection{Zerlegung und Radienklassen}

\begin{definition}[Fundamentale Torusstruktur]\label{def:t4}
Die fundamentale Ebene der FFGFT ist der kompakte 4-Torus
\[
  T^4 \;=\; S^1_x \times S^1_y \times S^1_z \times S^1_m ,
\]
mit drei \emph{r\"aumlichen} Kreisen $S^1_{x,y,z}$ und einem
\emph{Masse-Kreis} $S^1_m$. Alle vier Komponenten sind kompakt.
\end{definition}

\begin{remark}[Die vierte Richtung als Ersatz-Zeit]\label{bem:struktur}
Die vierte Torusrichtung $S^1_m$ ist nicht eine Massedimension \emph{neben}
einer separaten Zeit, sondern die \emph{Ersatz-Zeit-Richtung} selbst. Masse und
Zeit sind zwei \emph{Lesarten} ein- und derselben Dimension; die
Fundamentalrelation $\tilde{T}\cdot m = 1$ ist das \"Ubersetzungsw\"orterbuch:
\[
  \underbrace{S^1_m \text{ (kompakt)}}_{\text{Masse-Lesart, diskretes Spektrum}}
  \;\;\xleftrightarrow{\;\;\tilde{T}\cdot m = 1\;\;}\;\;
  \underbrace{\mathbb{R}_t \text{ (offen)}}_{\text{Zeit-Lesart, kontinuierlicher Fluss}}.
\]
Es gibt keine Zeit ``zus\"atzlich'' zu den vier Richtungen; die Zeit ist die
dekompakte Lesart der vierten Richtung (vgl.\ Dok.~185). Die Zwei-Typen-Struktur
h\"angt allein an der Trennung \emph{kleiner} und \emph{gro\ss er}
Kompaktifizierungsradien und ist von der Benennung unabh\"angig.
\end{remark}

Die vier Richtungen tragen zwei getrennte Radienklassen:
\begin{align}
  R_x = R_y = R_z &\;\sim\; R_H
  && \text{(kosmologischer Horizont, ``gro\ss '')},\\
  R_m &\;\sim\; \xi \cdot \ell_P
  && \text{(Planck-Skala, ``klein'')}.
\end{align}

\begin{remark}[Status der beiden Skalen]\label{bem:skalen}
$R_H$ ist hier die \emph{deklarierte externe} (importierte) kosmologische Skala,
\emph{nicht} aus $\xi$ hergeleitet. Intrinsisch ist allein die \emph{Form} -- die
Radienklassen-Trennung gro\ss/klein bzw.\ das dimensionslose Verh\"altnis
$R_H/\ell_P$ --, nicht der absolute \emph{Faktor} (Dok.~190, P20). Die
Zwei-Typen-Struktur dieses Dokuments h\"angt allein an dieser Klassen-Trennung,
nicht am Zahlenwert von $R_H$. Umgekehrt setzt $R_m\sim\xi\ell_P$ die \emph{Spitze}
des KK-Turms ($1/R_m\sim\xi^{-1}m_P$, transplanckisch); die beobachteten
Teilchenmassen liegen \"uber die Rekursion $\xi^n$ weiter unten auf der Leiter
(Dok.~159, 258/259).
\end{remark}

\subsection{Kaluza-Klein-Lesart}

Ein kompakter Kreis vom Radius $R$ tr\"agt einen diskreten Massenturm
$m_n \sim n/R$, $n \in \mathbb{Z}$. Die beiden Radienklassen erzeugen dadurch
qualitativ verschiedene Physik:

\begin{center}
\begin{tabular}{lll}
\hline
\textbf{Kreis} & \textbf{Radius} & \textbf{Modenturm} \\
\hline
$S^1_{x,y,z}$ & $R\sim R_H$ (gro\ss) & L\"ucke $\sim 1/R_H \approx 0$, quasi-kontinuierlich $\Rightarrow$ Raum \\
$S^1_m$ & $R\sim \xi\ell_P$ (klein) & L\"ucke riesig, diskret und schwer $\Rightarrow$ Massenspektrum \\
\hline
\end{tabular}
\end{center}

\begin{proposition}[$\tilde T m = 1$ als Kompaktifizierungsrelation]\label{prop:tm}
In nat\"urlichen Einheiten hat ein Kompaktifizierungsradius $R$ die Dimension
einer L\"ange bzw.\ inneren Zeit, $\tilde{T}\sim R$, w\"ahrend die erzeugte
Massenskala $m \sim 1/R$ ist. Die Fundamentalrelation
\[
  \boxed{\;\tilde{T}\cdot m = 1\;}
\]
ist damit strukturell die Kaluza-Klein-Relation f\"ur die Massedimension.
\end{proposition}

% ============================================================================
\section{Zwei Typen von Reduktion}
% ============================================================================

Die \emph{Abw\"arts}-Richtung (Reduktion von $T^4$) verl\"auft \"uber zwei
geometrische Operationen, die hier entwickelt werden. Die verlustfreie
\emph{Aufw\"arts}-Operation (Typ~III, repr\"asentationelle \"Ubersetzung) folgt
in Abschnitt~8.

\subsection{Typ~I: Dualit\"ats-Dekompaktifizierung (Masse $\to$ Zeit)}

\begin{definition}[Abrollen der Massedimension]\label{def:typeI}
Die Operation
\[
  S^1_m \;\xrightarrow{\;\text{\"Uberlagerung}\;}\; \mathbb{R}_t,
  \qquad p:\mathbb{R}_t \to S^1_m,\;\; t \mapsto e^{2\pi i\, t / R_m},
\]
rollt den kompakten Massekreis \"uber seine universelle \"Uberlagerung in die
nichtkompakte Zeitrichtung ab. Sie wird durch $\tilde{T}\cdot m = 1$ gesteuert.
\end{definition}

\begin{theorem}[Typ~I ist im Modeninhalt informationserhaltend]\label{satz:typeI}
Das Abrollen $S^1_m \to \mathbb{R}_t$ ist eine Einbettung, keine Projektion.
Periodische Strukturen auf $S^1_m$ heben sich eindeutig in $\mathbb{R}_t$; der
Funktionenraum auf $\mathbb{R}_t$ ist echt gr\"o\ss er. Es wird kein Modeninhalt
vergessen; vielmehr entsteht eine zus\"atzliche nichtkompakte Richtung -- die
Zeit.
\end{theorem}

\begin{beweis}
(Skizze) Die universelle \"Uberlagerung $p:\mathbb{R}\to S^1$ ist ein lokaler
Hom\"oomorphismus und surjektiv. Eine $R_m$-periodische Funktion auf $\mathbb{R}$
entspricht bijektiv einer Funktion auf $S^1_m$ (Fourier-Reihe, diskretes
Spektrum). Der umgekehrte Weg (Kompaktifizierung $\mathbb{R}\to S^1$) ist
verlustbehaftet, da nichtperiodische Anteile nicht herabsteigen. Das Abrollen
geht somit in die \emph{einbettende}, nicht in die projizierende Richtung.
\end{beweis}

\begin{corollary}[Zeit als Einbettungspreis]
Der ``Preis'' der Dekompaktifizierung ist kein Informationsverlust, sondern das
Auftreten der Zeitrichtung selbst (Dok.~185). Was nicht direkt sichtbar bleibt,
ist allein die \emph{Kompaktheit}: die Periodizit\"at von $S^1_m$ erscheint im
Beobachtbaren nicht als Geometrie, sondern indirekt -- als diskrete Struktur des
Spektrums.
\end{corollary}

\begin{remark}[Worauf sich ``nicht total verlustfrei'' bezieht]\label{bem:residual}
Es geht nicht um absoluten Verlust, sondern um \emph{lesart-abh\"angige
Sichtbarkeit}. In der Masse-Lesart ist die Kompaktheit manifest -- das Spektrum
ist diskret, die Identifikation $t \sim t + R_m$ direkt ausgedr\"uckt. In der
Zeit-Lesart ist dieselbe Kompaktheit nur \emph{latent}: aus kontinuierlicher
Zeit allein ist die topologische Identifikation nicht eindeutig
rekonstruierbar (die Windung bleibt mehrdeutig). Die Dualit\"at garantiert volle
\"Ubersetzbarkeit des Spektrums; die Identifikation tr\"agt die Zeit-Lesart nur
implizit -- \"uber die Diskretheit des Spektrums. Genau dort setzt $\xi^n$ an.
\end{remark}

\subsection{Typ~II: geometrische Projektion (r\"aumlich)}

\begin{definition}[Vergessen einer Raumrichtung]\label{def:typeII}
Eine Projektion $\pi:\mathbb{R}^{k}\to\mathbb{R}^{k-1}$ mit
$\ker(\pi)=\mathbb{R}$ vergisst eine r\"aumliche Richtung. Sie ist nicht
injektiv, also nicht umkehrbar.
\end{definition}

\begin{theorem}[Typ~II ist streng verlustbehaftet]\label{satz:typeII}
Bei einer r\"aumlichen Projektion $\mathbb{R}^k \to \mathbb{R}^{k-1}$ kollabiert
eine vollst\"andige Faser $\mathbb{R}$. Die Information \"uber die vergessene
Richtung ist im Bild nicht enthalten und ohne Zusatzdaten (Randbedingungen,
holographische Vollst\"andigkeit) nicht rekonstruierbar. Die Rekursion $\xi^n$
wirkt auf diesen Schritt \emph{nicht}.
\end{theorem}

\begin{center}
\begin{tabular}{lll}
\hline
 & \textbf{Typ~I (Masse$\to$Zeit)} & \textbf{Typ~II (r\"aumlich)} \\
\hline
Operation & Abrollen / \"Uberlagerung & Vergessen / Projektion \\
Richtung & einbettend & projizierend \\
Modeninhalt & erhalten & verloren \\
Reversibilit\"at & dual/spektral r\"uckgewinnbar & nicht reversibel \\
Rolle von $\xi^n$ & Spektralsignatur & wirkungslos \\
Residuum & Kompaktheit nur inferiert & echter Geometrieverlust \\
\hline
\end{tabular}
\end{center}

% ============================================================================
\section{Die formale Kette, Stufe f\"ur Stufe}
% ============================================================================

\subsection{Stufe~1: $D4 \to D3 + \text{Zeit}$ (Typ~I)}
\[
  T^4 \;\longrightarrow\; \underbrace{T^3}_{\to\,\mathbb{R}^3}
  \;\times\; \underbrace{\mathbb{R}_t}_{\text{aus } S^1_m}.
\]
Stufe~1 b\"undelt \emph{zwei} informationserhaltende Operationen \emph{verschiedenen}
Mechanismus: (Ia) der \emph{Massekreis} rollt gem\"a\ss\ Satz~\ref{satz:typeI} \"uber
die Dualit\"at $\tilde T m=1$ in $\mathbb{R}_t$ ab (Dualit\"ats-Dekompaktifizierung);
(Ib) die drei \emph{Raumkreise} dekompaktifizieren durch den \emph{Grenzprozess}
gro\ss er Radien ($R_{x,y,z}\to\infty$ bei $r\ll R_H$) zu $\mathbb{R}^3$
(Limes-Dekompaktifizierung). Beide sind einbettend, keine Projektion; (Ib) ist aber
\emph{kein} Dualit\"atsschritt, sondern ein Limes -- siehe die Pr\"azisierung in
Abschnitt~\ref{sec:offen}, Punkt~3.

\textit{Invariant:} $\tilde{T}\cdot m = 1$ (Prop.~\ref{prop:tm}); im r\"aumlichen
Limes die lokale Flachheit.\quad
\textit{Verlust:} keiner im Modeninhalt; nur Sichtbarkeit der Kompaktheit
(Bem.~\ref{bem:residual}).

\subsection{Stufe~2: $D3 \to D2$ (Typ~II, holographisch)}
\[
  \pi_2:\; \mathbb{R}^3 \;\longrightarrow\; \mathbb{R}^2,
  \qquad \ker(\pi_2) = \mathbb{R}\;\;(\text{Tiefenrichtung}).
\]
\textit{Invariant:} Fl\"ache $A$.\quad
\textit{Holographie:} Stufe~2 ist mathematisch das holographische Prinzip; im
Horizontfall der Bekenstein-Hawking-Bound $S \le A/(4\ell_P^2)$ (Anschluss
Dok.~254/268).\quad
\textit{Verlust:} echte Tiefenrichtung.

\subsection{Stufe~3: $D2 \to D1$ (Typ~II)}
\[
  \pi_3:\; \mathbb{R}^2 \;\longrightarrow\; \mathbb{R},
  \qquad \ker(\pi_3) = \mathbb{R}.
\]
\textit{Invariant:} L\"ange $L$.\quad
\textit{Verlust:} zweite Raumrichtung; kein Fl\"achenma\ss\ mehr.

\subsection{Stufe~4: $D1 \to D0$ (Typ~II, Totalverlust)}
\[
  \pi_4:\; \mathbb{R} \;\longrightarrow\; \{\mathrm{pt}\}.
\]
\textit{Invariant:} nur diskretes skalares Ma\ss\ ($n\in\mathbb{Z}$).\quad
\textit{Verlust:} vollst\"andige geometrische Struktur.

% ============================================================================
\section{Reversibilit\"at, pr\"azise pro Typ}
% ============================================================================

\begin{theorem}[Reversibilit\"at der Gesamtkette]\label{satz:rev}
Sei $\Pi = \pi_4\circ\pi_3\circ\pi_2$ die r\"aumliche Teilkette
$\mathbb{R}^3 \to \{\mathrm{pt}\}$ und $D$ die Typ-I-Operation
$S^1_m \to \mathbb{R}_t$. Dann gilt:
\begin{enumerate}
  \item[(i)] $D$ ist als Dualit\"at (Fourier/KK) im Modeninhalt umkehrbar; die
  Windung bleibt mehrdeutig. $D$ ist ``fast reversibel''.
  \item[(ii)] $\Pi$ ist nicht invertierbar: es gibt keine stetige Abbildung
  $\rho:\{\mathrm{pt}\}\to\mathbb{R}^3$ mit $\rho\circ\Pi=\mathrm{id}$.
\end{enumerate}
\end{theorem}

\begin{beweis}
(Skizze) (i) folgt aus Satz~\ref{satz:typeI}: das Abrollen ist einbettend, das
diskrete Spektrum kodiert den vollen periodischen Inhalt; nur die Wahl des Lifts
(Windung) ist nicht eindeutig. (ii) folgt aus Satz~\ref{satz:typeII}: jedes
$\pi_k$ hat $\ker = \mathbb{R}\neq\{0\}$. Eine Linksinverse m\"usste die
kollabierte Faserstruktur rekonstruieren; diese Information ist im Bild
$\{\mathrm{pt}\}$ nicht vorhanden.
\end{beweis}

\begin{corollary}[Aufl\"osung der Eingangsspannung]
Bild~(A) und Bild~(B) sind beide korrekt, betreffen aber verschiedene
Operationen. Die Nicht-Reversibilit\"at (A) gilt f\"ur die \emph{r\"aumlichen}
Projektionen (Typ~II). Die Dekompaktifizierung zur Zeitrichtung (B) ist der
\emph{einbettende} Typ-I-Schritt -- er erzeugt eine Dimension, statt eine zu
vergessen. Es liegt kein Widerspruch vor.
\end{corollary}

% ============================================================================
\section{Rolle der Rekursion $\xi^n$}
% ============================================================================

Die Rekursion $R:\xi^n \mapsto \xi^{n+1}$ ist \emph{keine} universelle
Kompensation f\"ur Projektionsverluste. Ihre Funktion ist eng begrenzt:
\begin{itemize}
  \item \textbf{Typ~I:} $\xi^n$ erzeugt das diskrete Massenspektrum
  (Teilchenleiter). Dieses Spektrum \emph{ist} der kompakte Modeninhalt von
  $S^1_m$, dual als Zeit-Frequenz-Struktur. $\xi^n$ ``f\"ullt'' nichts auf,
  sondern ist die Spektralsignatur, an der die verborgene Kompaktheit erkennbar
  bleibt (Bem.~\ref{bem:residual}).
  \item \textbf{Typ~II:} $\xi^n$ wirkt nicht. R\"aumliche Projektionsverluste
  sind echt und werden von der Rekursion nicht ber\"uhrt.
\end{itemize}

\begin{remark}[Korrektur einer naheliegenden Fehldarstellung]
Eine Darstellung, die der Rekursion eine \"uber alle vier Stufen monoton
abnehmende ``Kompensationskraft'' zuschreibt (voll/partiell/schwach/keine), ist
irref\"uhrend. Sie unterstellt eine gleichf\"ormige verlustbehaftete Leiter, die
es nicht gibt. Die Rekursion geh\"ort ausschlie\ss lich zum Typ-I-Schritt.
\end{remark}

% ============================================================================
\section{\"Ubersicht}
% ============================================================================

\begin{table}[h!]
\centering
\begin{tabular}{lllll}
\hline
\textbf{Stufe} & \textbf{Operation} & \textbf{Typ} & \textbf{Invariant} & \textbf{$\xi^n$} \\
\hline
$D4\to D3{+}t$ & $S^1_m\to\mathbb{R}_t$ (abrollen) & I & $\tilde{T}m=1$ & Spektralsignatur \\
$D3\to D2$ & Raum vergessen & II & Fl\"ache $A$ & -- \\
$D2\to D1$ & Raum vergessen & II & L\"ange $L$ & -- \\
$D1\to D0$ & Raum vergessen & II & $n\in\mathbb{Z}$ & -- \\
\hline
\end{tabular}
\caption{Die Kette $T^4 \to T^0$ nach Typ getrennt.}
\end{table}

\[
\underbrace{T^4 \;\longrightarrow\; \mathbb{R}^3\times\mathbb{R}_t}_{\text{Typ I: einbettend, fast reversibel}}
\;\longrightarrow\;
\underbrace{\mathbb{R}^2 \;\longrightarrow\; \mathbb{R} \;\longrightarrow\; \{\mathrm{pt}\}}_{\text{Typ II: projizierend, nicht reversibel}}
\]

% ============================================================================
\section{Die Aufw\"artsrichtung: repr\"asentationelle \"Ubersetzung (Typ~III)}
% ============================================================================

Die bisherige Kette beschreibt ausschlie\ss lich die \emph{Abw\"arts}-Richtung
(Reduktion). Es gibt eine Gegenrichtung -- aber sie ist \emph{nicht} geometrisch.

\begin{definition}[Typ~III: repr\"asentationelle \"Ubersetzung]\label{def:typeIII}
Die Operation
\[
  T^4 \;\xleftrightarrow{\;\text{bijektiv}\;}\; \mathcal{H}
\]
\"ubersetzt den FFGFT-Modenformalismus auf $T^4$ in den (unendlich\-dimensionalen)
Hilbertraum $\mathcal{H}$ bzw.\ die \"aquivalente Matrix-Algebra (Dok.~230, 206,
231). Sie ist keine Approximation, sondern eine \emph{bijektive
Strukturaussage}: jede QM-Rechnung ist in FFGFT-Sprache abbildbar und umgekehrt.
\end{definition}

\begin{important}[Die dimensionale Asymmetrie]
Abw\"arts (geometrisch) und aufw\"arts (repr\"asentationell) sind verschiedene
Arten von Operation mit \emph{entgegengesetzten} Informationseigenschaften:
\begin{itemize}
  \item \textbf{Abw\"arts (Typ~II, geometrisch):} verlustbehaftet, nicht
  reversibel -- eine Faser kollabiert (Satz~\ref{satz:typeII}).
  \item \textbf{Aufw\"arts (Typ~III, repr\"asentationell):} bijektiv,
  verlustfrei -- es werden \emph{keine} physikalischen Dimensionen hinzugef\"ugt,
  sondern derselbe 4D-Inhalt in einer h\"oherdimensionalen \emph{Sprache}
  ausgedr\"uckt, in der er ohne Reduktion Platz hat.
\end{itemize}
Die Verlustfreiheit folgt also \emph{gerade aus} dem repr\"asentationellen
Charakter -- nicht aus einem Gewinn echter Extra-Dimensionen.
\end{important}

\begin{remark}[Keine geometrisch h\"oheren Dimensionen]
Eine geometrische Erweiterung $T^4 \to T^5 \to \dots$ ist in der FFGFT
ausgeschlossen (Dok.~171): Komplexit\"at entsteht aus \emph{fraktaler Tiefe}
($D_f = 3-\xi$, Selbst\"ahnlichkeit), nicht aus \emph{mehr} Dimensionen. Dies
ist der bewusste Gegensatz zur Stringtheorie (10/11D, Calabi-Yau,
Landschaftsproblem). ``H\"oher'' bedeutet in der FFGFT daher stets Typ~III
(repr\"asentationell), nie geometrisch.
\end{remark}

Damit hat die FFGFT drei dimensions\"andernde Operationen:

\begin{center}
\begin{tabular}{llll}
\hline
\textbf{Typ} & \textbf{Operation} & \textbf{Charakter} & \textbf{Information} \\
\hline
I & Masse\,$\leftrightarrow$\,Zeit ($\tilde T m=1$) & geometrisch & fast reversibel \\
II & r\"aumliche Projektion (runter) & geometrisch & verlustbehaftet \\
III & $T^4\leftrightarrow\mathcal{H}$/Matrix (hoch) & repr\"asentationell & \textbf{bijektiv} \\
\hline
\end{tabular}
\end{center}

\noindent
Die Bijektivit\"at von Typ~III ist die spiegelbildliche Aussage zur
Nicht-Reversibilit\"at von Typ~II (Satz~\ref{satz:rev}): geometrisch verliert
man beim Reduzieren, repr\"asentationell verliert man beim \"Ubersetzen nichts.

% ============================================================================
\section{Physikalische Pr\"ufbarkeit der Identifikation $t\sim t+R_m$}
% ============================================================================

Die topologische Identifikation $t \sim t + R_m$ ist keine reine
Begleitstruktur, sondern eine physikalische Observable. In der Zeit-Lesart ist
ihre Periode \emph{exakt} die fundamentale Umlaufzeit
\[
  T_\text{rev} \;=\; \frac{R_m}{c} \;=\; \frac{\xi\,\ell_P}{c}
  \;\approx\; 7{,}2\times10^{-48}\,\text{s},
  \qquad
  \Delta\phi \;=\; 2\pi\xi \;\approx\; 8{,}4\times10^{-4}\ \text{rad/Umlauf},
\]
sodass $t\sim t+R_m$ identisch ist mit $t\sim t+c\,T_\text{rev}$ (Dok.~183/185).
Die Kompaktheit der vierten Richtung \emph{ist} die Revival-Zeit.

\subsection{Zwei pr\"ufbare Lesarten}
\begin{itemize}
  \item \textbf{Masse-Lesart (reif, quantitativ).} Das diskrete Modenspektrum
  erzeugt die Teilchenmassen als Torus-Obert\"one mit Schrittweite $1/D$; daraus
  folgen Koide, $g\!-\!2$ und drei Generationen (Dok.~159, 258/259). Hier ist
  die Kompaktheit zahlenm\"a\ss ig eingel\"ost.
  \item \textbf{Zeit-Lesart (falsifizierbar, derzeit unter Aufl\"osung).} Die
  Periodizit\"at erzeugt eine CHSH-Abweichung von der Tsirelson-Schranke der
  Gr\"o\ss enordnung $\xi$. F\"ur 73 Qubits sagt die FFGFT
  $\Delta \approx 5{,}4\times10^{-4}$ voraus (aufl\"osbar bei
  $\sigma\sim10^{-4}$). Der Hardwaretest IBM~Heron~r2 (Mai~2026) best\"atigt die
  Bell-Verletzung ($S=2{,}74$) und die Lauff\"ahigkeit der T0-Schaltkreise; der
  $\xi$-Effekt liegt jedoch unter dem NISQ-Rauschen und ist noch nicht
  aufl\"osbar (Dok.~147).
\end{itemize}

\subsection{Skalierung: linear versus logarithmisch}

\begin{important}[Regime-Abh\"angigkeit der Phasenakkumulation]
Eine naive \emph{koh\"arente Linearakkumulation} $r\cdot\Delta\phi$ \"uber $r$
Uml\"aufe erg\"abe physikalisch absurde Phasendriften (Gr\"o\ss enordnung
$10^{44}$~rad/s) und st\"unde im Widerspruch zu funktionierenden Atomuhren und
tiefen Quantenschaltkreisen. Die labor-pr\"ufbare Vorhersage skaliert daher
\emph{logarithmisch},
\[
  \delta\phi \;\sim\; \frac{\xi\,\ln(r)}{D_f}
  \;\approx\; 2\times10^{-4} \quad (D_f = 3-\xi),
\]
nicht linear. Die lineare Form $r\cdot\Delta\phi$ aus Dok.~185 ist die
\emph{Worst-Case-Schranke des RSA-Arguments} (Koh\"arenz \"uber $r\approx N$
Zyklen), nicht die Laborvorhersage. Diese Regime-Trennung erkl\"art, warum tiefe
koh\"arente Schaltkreise laufen, w\"ahrend die RSA-Faktorisierung gro\ss er
Schl\"ussel geometrisch gesch\"utzt bleibt.
\end{important}

\paragraph{Fazit.}
$t\sim t+R_m$ ist im FFGFT-Korpus als Observable gef\"uhrt ($T_\text{rev}$,
CHSH-Abweichung $\sim\xi$): in der Masse-Lesart quantitativ best\"atigt, in der
Zeit-Lesart falsifizierbar, aber derzeit knapp unter der experimentellen
Aufl\"osung. Keine unfalsifizierbare Begleitstruktur.

% ============================================================================
\section{Einordnung und offene Fragen}\label{sec:offen}
% ============================================================================

\paragraph{Anschl\"usse.}
Dok.~185 (Zeit als Einbettungspreis) liefert die ontologische Grundlage des
Typ-I-Schritts; Dok.~268 (holographische Projektion) und Dok.~254
(Fl\"achen-Schrankentheorem) den Rahmen f\"ur Stufe~2; Dok.~206
(Triangle--Matrix-Reduktion) best\"atigt, dass die 4D-Ebene generativ ist;
Dok.~147/183 liefern die zeit-seitige Pr\"ufbarkeit. Die Aufw\"arts-Operation
(Typ~III) ist in Dok.~230 (bijektive Hilbertraum-\"Ubersetzung), Dok.~231
(notwendige Erweiterungen) und Dok.~206 (Matrix-Darstellung) ausgearbeitet;
Dok.~171 begr\"undet den Ausschluss geometrisch h\"oherer Dimensionen.

\paragraph{Offene Punkte.}
\begin{enumerate}
  \item \textbf{Strukturannahme -- best\"atigt.} Die Zerlegung
  $T^4 = S^1_{x,y,z}\times S^1_m$ mit $S^1_m$ als Ersatz-Zeit-Richtung und
  $\tilde{T}\cdot m = 1$ als \"Ubersetzung ist festgelegt
  (Bem.~\ref{bem:struktur}).
  \item \textbf{Skalierung linear/logarithmisch.} Die lineare Worst-Case-Schranke
  $r\cdot\Delta\phi$ (Dok.~185, RSA-Regime) und die logarithmische
  Laborvorhersage $\xi\ln(r)/D_f$ (Dok.~147) sollten explizit als
  regime-abh\"angig gekennzeichnet werden, um Widerspruchsfreiheit zu sichern.
  \item \textbf{Spatialer Grenzfall -- pr\"azisiert.} Die Dekompaktifizierung
  $S^1_{x,y,z}\to\mathbb{R}^3$ ist ein \emph{Grenzprozess} ($R_{x,y,z}\to\infty$),
  weder die Typ-I-Dualit\"at noch eine Typ-II-Projektion, sondern eine eigene
  \emph{limes-artige} informationserhaltende Operation~(Ib). Sie wird daher in
  Stufe~1 ausdr\"ucklich von der Massekreis-Dualit\"at~(Ia) getrennt gef\"uhrt;
  beide teilen den einbettenden Charakter, nicht den Mechanismus.
  \item \textbf{Holographische Vollst\"andigkeit.} Inwieweit l\"asst sich der
  Typ-II-Verlust bei $D3\to D2$ \"uber Randdaten exakt (nicht nur teilweise)
  rekonstruieren? Anschluss Dok.~254/268.
\end{enumerate}
